\documentclass[10pt,a4paper]{amsart}
\usepackage[margin=19mm]{geometry}
\usepackage{amsmath,amssymb,mathtools}
\usepackage[scr=rsfs,cal=euler]{mathalfa}
\usepackage{xfrac}
\usepackage[unicode,hidelinks,bookmarks=false]{hyperref}
\newcommand{\ie}{i.e.\ }
\newcommand{\cf}{cf.\ }
\newcommand{\resp}{respectively\ }
\newcommand{\Subsection}{\subsection}
\providecommand{\nobreakdash}{\nobreak\hbox{-}}
\setlength{\emergencystretch}{2em}
\multlinegap0pt
\usepackage{bm}
\usepackage{bbm}
\usepackage{dsfont}
\usepackage{xspace}
\usepackage{cancel}
\usepackage{mathtools}
\usepackage[shortlabels]{enumitem}
\usepackage{amsthm}
\makeatletter
\@ifundefined{thm}{\newtheorem{thm}{Theorem}}{}
\@ifundefined{prop}{\newtheorem{prop}[thm]{Proposition}}{}
\makeatother
\title[Differential inclusions and quasi-Lyapunov functions]{Differential inclusions and quasi-Lyapunov functions}
\author{Martin Ivanov, Mikhail Krastanov, Nadezhda Ribarska}
\date{Source: arXiv:2511.06100v1 (CC BY 4.0)}
\begin{document}
\maketitle

\noindent\emph{Source and licence.} Differential inclusions and quasi-Lyapunov functions. Version arXiv:2511.06100v1 (CC BY 4.0), distributed under the Creative Commons Attribution 4.0 International licence, as recorded in the arXiv licence metadata for this exact version and shown on the abstract page https://arxiv.org/abs/2511.06100v1. Full rights notice: licenses/arxiv-2511.06100-CC-BY-4.0.txt.\par\smallskip

\noindent\textit{Editorial scope.} Counted unit: the Thm whose source label is \texttt{sol} in the article (source0.tex lines 292--300, source label \texttt{sol}), delivered together with the article's own complete original proof of it (source0.tex lines 301--317, one \texttt{proof} environment of 17 lines). No mathematics is written, repaired or re-derived: statement and proof are the article's own text line for line. Required context retained uncounted: ulimit-sol (lines 185--191); esolution (lines 221--226); time (lines 269--271); time2 (lines 277--281). Every definition, display and label of those lines is retained unchanged. Omitted from this package and not redistributed: 1--184, 192--220, 227--268, 272--276, 282--291, 318--621. Nothing inside a retained excerpt is omitted or altered: every retained body is delivered exactly as the article's lines are written, so no omission receipt of any kind accompanies this package. Citations: the retained excerpts carry no citation command at all, so no bibliography entry of the article is transcribed and no reference is redistributed. Reference adaptation: none was needed and none was made; every reference inside the delivered unit resolves inside it, so no unresolved reference or citation is produced. Printed slips kept literal: no printed word, symbol, hypothesis or number of the counted statement or of its proof was altered. Layout and class: the article's own class is \texttt{article} with packages including babel, inputenc, amsmath, amssymb, amsthm, amscd, amsfonts, newlfont, enumerate, graphicx, ifsym, titlefoot, hyperref; the wrapper is the lane's pinned \texttt{amsart} editorial wrapper at 10pt on A4, which restates the theorem-like environments the retained text needs and adds the article's own macro definitions that the retained text needs (none), with extra wrapper packages (bm, bbm, dsfont, xspace, cancel, mathtools, enumitem). The delivered numbering is the unit's own. Assets: no retained span contains a figure environment, a graphics-inclusion command, a PDF-page inclusion, a table or a TikZ picture; no image file is copied into this package, the package declares no asset dependency and no image file is redistributed with it. Licence: version arXiv:2511.06100v1 of this article is distributed under the Creative Commons Attribution 4.0 International licence, as recorded in the arXiv licence metadata for this exact version and shown on the abstract page https://arxiv.org/abs/2511.06100v1. A full rights notice is delivered at licenses/arxiv-2511.06100-CC-BY-4.0.txt. Characters: every retained excerpt is pure ASCII, so the delivered engine, XeLaTeX with its default Latin Modern text font, reports no missing character.
\begin{prop}\label{ulimit-sol}
    Let $D$ be a locally closed subset of $\mathbb{R}^n$ and $F: \overline{D} \rightrightarrows \mathbb{R}^n$ be uniformly bounded. Let $\mathcal{U}$ be a $\sigma$-relatively open partition of $D$ and $(\mathcal{U},\mathcal{G})$ be an invariant partial $\varepsilon$-approximation of $F$. Then:
    \begin{enumerate}[(i)]
        \item for each absolutely continuous function $\varphi: [0,T] \to \cup\left(\mathbf{Green}\,\mathcal{U}\right)$ with respect to which the elements of $\mathbf{Green}\,\mathcal{U}$ are forward invariant, the function $$t \mapsto \alpha(t), \text{where } \varphi(t) \in U_{\alpha(t)}$$ is monotone increasing;
        \item if $\varphi_k : [0,T] \to \cup\left(\mathbf{Green}\,\mathcal{U}\right)$, $k=1,2,\dots$ is a sequence of $\varepsilon$-solutions of $\dot{x} \in F(x)$ with respect to which the elements of $\mathbf{Green}\,\mathcal{U}$ are forward invariant and $\varphi : [0,T] \to \cup\left(\mathbf{Green}\,\mathcal{U}\right)$ is its uniform limit, then $\varphi$ is an $\varepsilon$-solution of $\dot{x} \in F(x)$.
    \end{enumerate}
\end{prop}

\begin{prop}\label{esolution}
    Let $D$ be an intersection of an open subset $S$ and a closed subset $K$ of $\mathbb{R}^n$ and $F : \overline{D} \rightrightarrows \mathbb{R}^n$ be a uniformly bounded multi-valued mapping with nonempty values.
    Let $\mathcal{U}$ be a $\sigma$-relatively open partition of $D$ and $(\mathcal{U},\mathcal{G})$ be an invariant partial $\varepsilon$-approximation of $F$. Let $w: D \to \mathbb{R}$ be a quasi-Lyapunov function for $(\mathcal{U},\mathcal{G})$.

    Then for each point $x_0 \in \overline{\cup(\mathbf{Green} \,\mathcal{U})} \cap D$ there exists an $\varepsilon$-solution $\varphi(\cdot)$ of $\dot{x} \in F(x)$ starting from $x_0$ and defined on the interval $[0,T)$ ($T=+\infty$ or $\varphi(t)$ tends to a point in $\mathbb{R}^n \setminus S$ whenever $t \to T$) such that $w(\varphi(t)) \ge w(x_0) + t$ for all $t \in [0,T)$. Moreover, if the starting point $x_0$ belongs to $\cup(\mathbf{Green} \,\mathcal{U})$, then each element of $\mathbf{Green} \,\mathcal{U}$ is forward invariant with respect to the corresponding $\varepsilon$-solution $\varphi(\cdot)$ starting from $x_0$.
\end{prop}

    \begin{equation}\label{time}
        T_0 \ge \frac{\mathrm{dist}(x_0,\mathbb{R}^n \setminus S)}{c+1} > 0.
    \end{equation}

    \begin{equation}\label{time2}
        \begin{aligned}
            T_k   {\ge} \frac{\mathrm{dist}(x_k,\mathbb{R}^n \setminus S)}{c+1} &\ge \frac{\mathrm{dist}(x_0,\mathbb{R}^n \setminus S) - \|x_k-x_0\|}{c+1}\\ &\ge T:= \frac{\mathrm{dist}(x_0,\mathbb{R}^n \setminus S)}{2(c+1)} > 0.
        \end{aligned}
    \end{equation}

\begin{thm}\label{sol}
    Let $D$ be a locally closed subset of $\mathbb{R}^n$ and $F : \overline{D} \rightrightarrows \mathbb{R}^n$ be a uniformly bounded multi-valued mapping with nonempty closed values. Let $\mathcal{U}^k$ be a $\sigma$-relatively open partition of $D$, $(\mathcal{U}^k,\mathcal{G}^k)$ be an invariant partial $\frac{1}{k}$-approximation of $F$ and $(\mathcal{U}^{k+1},\mathcal{G}^{k+1})$ be a refinement of $(\mathcal{U}^k,\mathcal{G}^k)$ for all $k=1,2,\dots$. We assume that for each approximation $(\mathcal{U}^k,\mathcal{G}^k)$ there exists a quasi-Lyapunov function $w_k : D \to \mathbb{R}$ such that $w_{k_1} \ge w_{k_2}$ on $D$ for every $k_1 \le k_2$.

    Then for each point $x_0 \in \bigcap_{k=1}^{\infty}\left(\overline{\bigcup(\mathbf{Green} \,\mathcal{U}^k)}\right) \bigcap D$ there exist $T>0$ and an absolutely continuous function $x : [0,T] \to D$ such that
    $$\left| \begin{array}{l}
    \dot{x}(t) \in F(x(t)) \text{ a.e. on } [0,T]\\
    x(0) = x_0\\
    \end{array} \right..$$
\end{thm}

\begin{proof}
    Fix an arbitrary point $x_0 \in \bigcap_{k=1}^{\infty}\left(\overline{\bigcup(\mathbf{Green} \,\mathcal{U}^k)}\right) \bigcap D$ and a positive integer $k$. Then $x_0 \in \left(\overline{\bigcup(\mathbf{Green} \,\mathcal{U}^k)} \cap D\right)$, so there exists $x_k \in \bigcup(\mathbf{Green} \,\mathcal{U}^k)$ such that $\|x_k-x_0\| < \frac{1}{k}$.

    Then, by applying Proposition \ref{esolution}, there exists a $\frac{1}{k}$-solution $\varphi_k(\cdot)$ of $\dot{x} \in F(x)$ starting from $x_k$ and defined on some maximal interval $[0,T_k)$ such that each element of $\mathbf{Green} \,\mathcal{U}^k$ is forward invariant with respect to $\varphi_k$ and $w_k(\varphi_k(t)) \ge w_k(x_k) + t$ for all $t \in [0,T_k)$.

    Let $D = S \cap K$, where $S$ is an open subset and $K$ is a closed subset of $\mathbb{R}^n$. According to Proposition \ref{esolution}, $T_k = +\infty$ or $\varphi_k(t)$ tends to a point in $\mathbb{R}^n \setminus S$ whenever $t \to T_k$. According to $(\ref{time})$, $$T_k \ge \frac{\mathrm{dist}(x_k,\mathbb{R}^n \setminus S)}{c+1} > 0.$$

    Repeating this for each positive integer $k$, we obtain a sequence $\{x_k\}_{k=1}^{\infty}$ which tends to $x_0$. Thus, there exists a constant $k_0 \in \mathbb{N}$ such that for all $k \ge k_0$ we have $\|x_k-x_0\| \le \frac{1}{2} \mathrm{dist}(x_0,\mathbb{R}^n \setminus S)$. Using the same arguments as in $(\ref{time2})$, for each $k \ge k_0$ we get $$T_k \ge T:= \frac{\mathrm{dist}(x_0,\mathbb{R}^n \setminus S)}{2(c+1)} > 0.$$

    Without loss of generality, we may think that the sequence $\{\varphi_k(\cdot)\}_{k=k_0}^{\infty}$ is uniformly convergent on the interval $[0,T]$ to an absolutely continuous function $\varphi : [0,T] \to D$.

    Fix an arbitrary positive integer $m \ge k_0$. Since $(\mathcal{U}^k,\mathcal{G}^k)$ is a refinement of $(\mathcal{U}^m,\mathcal{G}^m)$ for each $k \ge m$, the elements of $\mathbf{Green} \,\mathcal{U}^m$ are forward invariant with respect to the trajectory $\varphi_k(\cdot)$ for each $k \ge m$. Since $w_m(\varphi_k(t)) \ge w_k(\varphi_k(t)) \ge w_k(x_k) + t \ge t$ for all $t \in [0,T]$, by taking the limit as $k \to +\infty$, we obtain that $w_m(\varphi(t)) \ge t$ for all $t \in [0,T]$. Thus, $\varphi(t) \in \cup(\mathbf{Green} \,\mathcal{U}^m)$ for all $t \in (0,T]$. According to Proposition \ref{ulimit-sol}(ii), $\varphi(\cdot)$ is a $\frac{1}{m}$-solution of $\dot{x} \in F(x)$, i.e. $$\dot{\varphi}(t) \in F(\varphi(t)) + \frac{1}{m} \overline{\mathbf{B}} \text{ a.e. on } [0,T].$$

    Since $\dot{\varphi}(t)$ does not depend on $m$, the multi-valued mapping $F$ is closed-valued and the union of countably many sets of measure zero is a set of measure zero, we obtain that $$\dot{\varphi}(t) \in F(\varphi(t)) \text{ a.e. on } [0,T].$$

    This completes the proof.
\end{proof}

\end{document}
